\section*{Ethical Considerations}
This work studies how model compression via layer pruning and continued
pretraining affects evaluated capability; it does not add a safety mechanism and
changes none of the underlying model's risks such as hallucination, biased
generation, or unsafe completion. The observed PPL--MMLU dissociation carries a
practical caution: a compressed model that looks healthy on aggregate
language-modelling metrics may still perform poorly on a target benchmark, and
deployments should evaluate the capabilities they intend to preserve rather than
rely on perplexity alone.

We collect no new human-subject data and recruit no annotators. All experiments
use previously released artifacts---the OLMo-2 base model, the Dolmino corpus, and
standard public benchmarks---under their respective licenses. Training and
large-scale evaluation consume energy. We reuse a common training recipe and
report checkpoint step explicitly rather than presenting unequal-step results as
compute matched. Any release should preserve the licenses and usage restrictions
of the underlying model, corpus, and benchmarks. An anonymous artifact should
include training/evaluation code, exact configs, environment locks, checkpoint
hashes, per-item prediction files where redistribution permits, and the actual
data-order/resume behavior; upstream model weights and benchmark text should
remain with their original providers.
